% Copyright (c) 2026 Marcio Luiz Lima de Oliveira and V. el Hakim.
% Distribution and modification are subject to LPPL version 1.3c.
% SPDX-License-Identifier: LPPL-1.3c
% Part of the LPPL-covered work listed in MANIFEST.txt.
% See LICENSE and LICENSING.txt for the licence and scope.

\section{Preserve the university front matter}
The formal pages follow Annex 3 of the University of Twente Doctoral Regulations
effective 1 January 2026~\cite{ut2026}. The model uses an uppercase title,
``DOCTORAL THESIS'', an approval page and a manuscript committee page; it omits
birth details. The rector's name follows that model. Confirm the final wording,
names and pages with the Doctorate Board Support Office before printing.

Article 13 requires Dutch and English titles and abstracts; both are provided
in \texttt{Abstract.tex}. Article 19 requires title-page approval before copies
are made. The 2026 Doctorate workflow includes this approval~\cite{utDoctorate}.

\texttt{Title\_Preamb.tex} keeps the cover, half-title, blank reverse, formal
title page, approval/colophon, committee and optional dedication together.
The approval page follows the concrete Annex 3 model, including the credits
and identifiers shown there. Article 15 describes a narrower set of permitted
reverse-page material; have the Support Office resolve any requested variation.

Change the metadata and committee entries, not the formal degree wording to
match a decorative theme. Use the actual number of supervisors and the correct
singular or plural role headings. The example committee is not a prescription
for committee membership. Joint doctorates need the applicable agreed model.

\section{Separate presentation from required content}
The teal palette, sans-serif body, chapter tabs, 170 by 240\,mm trim size and
optional dedication are this layout's choices. They are not presented as
university requirements. \autoref{chap:06} explains how to change them.
Funding and affiliation examples are fictional; replace them with accurate
statements. The UT logo is retained as a separate asset, not redrawn or
recoloured. Institutional identity rules apply to its use.
